%%%%%%%%%%%%%%%%%%%%%%%%%%%%%%%%%
%
%       EXERCÍCIO
%
%%%%%%%%%%%%%%%%%%%%%%%%%%%%%%%%%

\ifdefstring{\atividade}{prova}{%
    \renewcommand{\valorquestao}{\ValorQ\ pontos}
}%

\begin{exercicioBanco}[\valorquestao]
    O tempo \(T\), em minutos, necessário para um operário processar certa peça é uma variável aleatória com a seguinte distribuição de probabilidade:
    \[
        \begin{array}{c|c|c|c|c|c|c}
            t & 2 & 3 & 4 & 5 & 6 & 7 \\
            \hline
            p(t) & 0,1 & 0,1 & 0,3 & 0,2 & 0,2 & 0,1
        \end{array}
    \]
    \begin{enumerate}[(a)]
        \item Calcule o tempo médio de processamento.
        %
        \item Cada peça processada paga ao operador R\$2,00 mas, se ele processa a peça em menos de 6 minutos, ganha R\$0,50 por minuto poupado. Por exemplo, se ele processa a peça em 4 minutos, ganha um bônus de R\$1,00. Encontre a distribuição, a média e a variância da variável aleatória \(G\): quantia paga por peça.
    \end{enumerate}
\end{exercicioBanco}

